\section{Cross-Dataset Transfer of the Certificate}
\label{sec:cross}
The certificate is a function of the estimated subspace $\bm M_s$ and the selection $S$; a
fair question is whether it is an artifact of one cohort. We re-estimate $\bm M_s$ under the
\emph{same} unified truncated-SVD numerics ($\varrho{=}10^{-2}$) on the independent
Chapman-Shaoxing-Ningbo 12-lead database \cite{zheng2020chapman} (different country, devices,
sampling; $\CrossNchap$ records) and compare to PTB-XL by principal angles and by the
per-configuration recoverability (rank + global $\kappa$ + per-lead $\eta$). We deliberately do
\emph{not} claim ``population independence''; we report what transfers.

\textbf{QRS subspace is transfer-stable.} For QRS the two datasets' rank-3 subspaces are nearly
aligned (max principal angle $\CrossQRSangMax^\circ$). The per-configuration conditioning
transfers for the well-posed sets: $\kappa$ for $\{$I,II,V2$\}$ is $\CrossKapTriplePtb$ (PTB-XL)
vs.\ $\CrossKapTripleChap$ (Chapman) and for $\{$I,II,V1,V3,V5$\}$ $\CrossKapSpanPtb$ vs.\
$\CrossKapSpanChap$; the near-rank-deficient $\{$V1,V2,V3$\}$ transfers less cleanly
($\CrossKapVvvPtb$ vs.\ $\CrossKapVvvChap$), as expected for an ill-conditioned set.

\textbf{limb-6 is effective rank $\CrossLimbRankPtb$ on both cohorts.} Under the unified
truncation the six limb leads span the frontal plane, so $\bm M_{s,S}$ is rank
$\CrossLimbRankPtb$ (the transverse singular value is truncated) on both PTB-XL and Chapman---we
do \emph{not} invert a near-zero singular value into a spurious $10^{5}$-scale $\kappa$.
Precordial unrecoverability is instead flagged, consistently across cohorts, by $\eta{>}0$
(e.g.\ V2 $\eta{=}\CrossLimbEtaVtwoPtb$ on PTB-XL vs.\ $\CrossLimbEtaVtwoChap$ on Chapman). So
the qualitative identifiable/unidentifiable verdict and the well-posed $\kappa$ ordering are
cohort-stable.

\textbf{ST/T third direction does not transfer.} For the low-amplitude segments the third
subspace direction differs sharply between datasets (max principal angle
$\CrossSTang^\circ$ for ST, $\CrossTang^\circ$ for T). We therefore restrict the transfer
claim: the QRS map and the limb-6 verdict are cohort-stable; the ST/T third spatial direction is
cohort-specific and its map should be re-estimated per dataset---the honest scope of transfer,
not a blanket invariance.

\textbf{Sensitivity to duplicated limb leads (\texttt{results/lead\_weighting.json}).} The six
limb leads are exact combinations of I and II, so a 12-lead PCA up-weights frontal-plane content.
Refitting $\bm M_s$ on the 8 algebraically-independent leads and lifting back via the fixed
transform $\bm T$, the \emph{identifiability verdict is robust}: for the spanning sets the
per-lead $\eta$ is unchanged (max $|\Delta\eta|{=}0$; all targets remain exactly identifiable),
and limb-6 precordial leads stay unidentifiable under both fits. The \emph{magnitudes} of the
low-amplitude ST/T subspace are, however, sensitive: the 12-lead and 8-lead ST/T subspaces
differ by large principal angles ($\LWSTang^\circ$/$\LWTang^\circ$) and $\kappa$ can shift
several-fold (e.g.\ ST $\{$I,II,V2$\}$: $\LWKapTwelve\!\to\!\LWKapEight$). We therefore read the QRS map quantitatively and the
ST/T \emph{graded} values qualitatively---the binary identifiable/unidentifiable verdict is
robust; the exact ST/T $\eta$/$\kappa$/ambiguity magnitudes are weighting-dependent.